% !PRINTED 106
\weberpage{106}
%<<CONT
\[ \mathfrak{A}_3)\quad \eta = i\,\frac{\sqrt{3}-1}{\sqrt{2}} \]
\[ \mathfrak{A}_4)\quad \eta = \frac{1+\sqrt{3}}{\sqrt{2}} \]
\[ \mathfrak{A}_5)\quad \eta = \frac{1+i\sqrt{3}}{\sqrt{2}}. \]

Es soll nun zuletzt noch der Fall einer einzigen Invariante
$i_1 = 4$ betrachtet werden. Die Bedingungen für $m$ bestehen
dann einfach darin, dass $m$ durch 4, 8, 16, aber keine höhere
Potenz von 2 theilbar sein kann, dass ungerade Primzahlen nur
einfach in $m$ aufgehen können, und dass mindestens einer der
Indexmoduln durch 4 theilbar sein muss. Die ersten Werthe
sind also
\[ m = 5,\ 13,\ 15,\ 16,\ 17,\ 20,\ 29,\ 33,\ 37,\ 39,\ 40\ \ldots \]

Wir wollen annehmen, es seien von den Indexmoduln
$c_1, c_2, \ldots c_\varrho$ durch 4 theilbar, $c_{\varrho+1}, c_{\varrho+2}, \ldots c_\mu$ durch 2, aber
nicht durch 4 theilbar. Es muss dann $\varrho$ mindestens gleich 1
sein, und folglich ist
\[ \begin{aligned}
&\delta_{1,1} = \delta_{1,2} = \cdots = \delta_{1,\varrho} = 4, \quad \delta_{1,\varrho+1} = \cdots = \delta_{1,\mu} = 2\\
&i_{1,1} = i_{1,2} = \cdots = i_{1,\varrho} = 1, \quad i_{1,\varrho+1} = \cdots = i_{1,\mu} = 2\\
&c_{1,1} = \frac{c_1}{4}, \cdots\cdots, c_{1,\varrho} = \frac{c_\varrho}{4}, \quad c_{1,\varrho+1} = \frac{c_{\varrho+1}}{2}, \cdots, c_{1,\mu} = \frac{c_\mu}{2}.
\end{aligned} \]

Nun sind die $e_{1,h}$ so zu bestimmen, dass von den Zahlen
\[ e_{1,1}, e_{1,2}, \ldots e_{1,\varrho} \]
wenigstens eine ungerade ist, und dass (wegen \S.\,21, 6.) von den
Zahlen
\[ e_{1,1}, e_{1,2}, \ldots e_{1,\varrho} \]
keine durch 4, und von den
\[ e_{1,\varrho+1}, e_{1,\varrho+2}, \ldots e_{1,\mu} \]
keine durch 2 theilbar sei. Eine Ausnahme bildet hierbei wieder
im Falle, wo $m$ durch 8 oder durch 16 theilbar ist, die Zahl $e_{1,h}$,
der keiner der Primfactoren von $m$ entspricht, die auch gerade
sein kann.

Dann erhalten wir für die Indices $\alpha$ der Zahlen $a$ die Congruenz:
\[ e_{1,1}\alpha_1 + \cdots + e_{1,\varrho}\alpha_\varrho + 2\,(e_{1,\varrho+1}\alpha_{\varrho+1} + \cdots + e_{1,\mu}\alpha_\mu) \equiv 0\ (\mathrm{mod}\ 4). \]
